% TOP_PROBLEM: {"id":30007730,"problem_number":"LOCAL-30007730","title":"Mental-health poverty feedbacks","category_id":21,"category":{"id":21,"name":"economics","display_name":"Economics","description":"Open economic theory statements, published research agendas, and proposed scoped specifications; question provenance and historical priority are qualified per record.","slug":"economics","order_index":21,"created_at":"2026-10-08T00:00:00Z"},"status":"research_question","external_url":null,"legacy_ids":[],"published":false,"statement":"Test five-year income and mental-health interactions between a cash grant and psychotherapy, and assess whether a specified dynamic model supports a psychological poverty trap.","economics":{"original_id":219,"field":"Health economics","kind":"empirical","formalization_basis":"adapted_research_question","source_status":"explicit_open","priority_status":"earliest_found","priority_note":"Earliest explicit statement located in this audit; no exhaustive priority search or first-ever claim. The mathematical specification below is adapted, not quoted from the source.","earliest_verified_reference":{"authors":["Matthew Ridley","Gautam Rao","Frank Schilbach","Vikram Patel"],"title":"Poverty, depression, and anxiety: Causal evidence and mechanisms","year":2020,"url":"https://economics.mit.edu/sites/default/files/inline-files/poverty-depression-anxiety-science.pdf","doi":"10.1126/science.aay0214","locator":"Review Summary, Outlook; full article p. 9, Poverty traps and Box 4","evidence_summary":"The authors explicitly prioritize testing mental-health poverty traps and evaluating combinations of economic and psychological support over longer horizons."},"current_reference":{"authors":["Matthew Ridley","Gautam Rao","Frank Schilbach","Vikram Patel"],"title":"Poverty, depression, and anxiety: Causal evidence and mechanisms","year":2020,"url":"https://economics.mit.edu/sites/default/files/inline-files/poverty-depression-anxiety-science.pdf","doi":"10.1126/science.aay0214","locator":"Review Summary, Outlook; full article p. 9, Poverty traps and Box 4","evidence_summary":"The authors explicitly prioritize testing mental-health poverty traps and evaluating combinations of economic and psychological support over longer horizons."},"formalization_note":"The source explicitly identifies the underlying gap or agenda. Population, horizon, interventions, estimands, and auxiliary assumptions are proposed operational choices, not a canonical source theorem. Partial later answers are compatible with this scoped question; the citation alone does not prove absence of every subsequent solution."},"statusQualification":"Published question or research agenda verified at the cited locator. The mathematical specification is adapted for this catalogue and includes additional assumptions; source publication does not establish first-ever priority or certify this exact model remains unsolved. The source explicitly identifies the underlying gap or agenda. Population, horizon, interventions, estimands, and auxiliary assumptions are proposed operational choices, not a canonical source theorem. Partial later answers are compatible with this scoped question; the citation alone does not prove absence of every subsequent solution.","statusReviewedAt":"2026-10-08"}
% FIRST_POSTED: 2026-10-08
% ENTRY_KIND: statement_only
% !TeX program = lualatex
\documentclass[11pt]{article}
\usepackage{fontspec}
\usepackage[a4paper,margin=25mm]{geometry}
\usepackage{amsmath,amssymb,mathtools}
\usepackage{xurl}
\usepackage[hidelinks]{hyperref}
\newcommand{\sref}[2]{\hyperlink{source:#1:#2}{[S#2]}}
\setlength{\emergencystretch}{3em}
\begin{document}
% problemId:30007730
\section*{Mental-health poverty feedbacks}\label{30007730}
\subsection{Definitions and mathematical statement}
Consider low-income adults screening positive for depression in one district of India. Randomize one-time treatment offers \(z=(c,m)\in\{0,1\}^2\), where \(c\) offers a defined cash grant and \(m\) offers a defined course of evidence-based psychotherapy. For adult \(i\), let \(y_{it}(z)\) be real income and \(s_{it}(z)\) a validated symptom score at years \(t=1,3,5\). Define factorial interactions
\[\gamma_t^y=E[y_{it}(1,1)-y_{it}(1,0)-y_{it}(0,1)+y_{it}(0,0)],\]
and define \(\gamma_t^s\) by replacing income with symptoms. Estimate these interactions and all four treatment means under randomized offers, consistency, and measured village spillovers. A lasting positive income effect and lower symptoms may indicate reinforcing gains, but do not alone establish multiple stable equilibria. To evaluate a psychological poverty trap, specify a dynamic state-transition model and test whether its nonlinearities predict persistent divergence after temporary interventions, beyond baseline heterogeneity and regression to the mean. State attrition assumptions and collect consumption and work outcomes. The source explicitly prioritizes testing mental-health poverty traps and effective combinations; the population, factorial intervention, and finite horizons are our operationalization.

\par\textbf{Formulation and scope.} The source explicitly identifies the underlying gap or agenda. Population, horizon, interventions, estimands, and auxiliary assumptions are proposed operational choices, not a canonical source theorem. Partial later answers are compatible with this scoped question; the citation alone does not prove absence of every subsequent solution.

\par\textbf{Open-question provenance.} An explicit question or research agenda was verified at the source locator. This does not by itself certify that every later version remains unresolved \sref{30007730}{1}.

\par\textbf{Historical priority.} Earliest explicit statement located in this audit; no exhaustive priority search or first-ever claim. The mathematical specification below is adapted, not quoted from the source.

\subsection{Short English statement}
Test five-year income and mental-health interactions between a cash grant and psychotherapy, and assess whether a specified dynamic model supports a psychological poverty trap.

\subsection{Sources}
\begin{itemize}\raggedright
\item[\textnormal{[S1]}]\hypertarget{source:30007730:1}{} Matthew Ridley, Gautam Rao, Frank Schilbach, Vikram Patel. 2020. Poverty, depression, and anxiety: Causal evidence and mechanisms. Review Summary, Outlook; full article p. 9, Poverty traps and Box 4.
\url{https://economics.mit.edu/sites/default/files/inline-files/poverty-depression-anxiety-science.pdf}.
DOI: 10.1126/science.aay0214.
{\small\textit{Earliest verified question/agenda reference; historical priority qualified below. The authors explicitly prioritize testing mental-health poverty traps and evaluating combinations of economic and psychological support over longer horizons.}}
\end{itemize}
\end{document}
